\documentclass[
landscape,
  a4paper,
  12pt,
  english,
  brazilian,
]{article}

\usepackage[]{fatec-article}
\usepackage{setspace}

\begin{document}

\section*{Instruções para o preenchimento}
\doublespacing
\begin{enumerate}
    \item O Diário de Bordo é usado para registrar atividades, progressos, ideias e desafios enfrentados em um projeto ou durante a rotina de trabalho. Serve como um registro cronológico e detalhado das operações diárias, facilitando a organização e o acompanhamento das tarefas.
    \doublespacing
    \item Durante o registro das atividades deve-se incluir detalhes como datas, horários, descrições de tarefas, nomes de participantes e observações relevantes.  Esta documentação contínua ajuda na avaliação do progresso de projetos ou atividades, permitindo ajustes e melhorias contínuas nos processos.
    \doublespacing
    \item Para evidenciar a realização das tarefas, você poderá utilizar a criação de anexos para adicionar anotações, fotos, prints, questionários, entre outros.
\end{enumerate}

 \begin{table}[]
\centering
 

\begin{tabular}{|>{\raggedright\arraybackslash}p{0.2\linewidth}|l|l|>{\raggedright\arraybackslash}p{0.2\linewidth}|>{\raggedright\arraybackslash}p{0.2\linewidth}|}
\hline
Nome da Atividade & Data de início & Data de término & Responsável pela atividade & Descrição da atividade realizada \\\hline
                  Início dos trabalhos para o projeto&                            31/07/2026&                31/07/2026&                  Tiago Rodrigues; Davi Almeida; Leandro Augusto&                                  Início dos planejamentos e realizações para o projeto no 6° Semestre.\\ \hline
                  Apresentação na integração&                            31/07/2026&                31/07/2026&                 Tiago Rodrigues; Davi Almeida; Leandro Augusto;&                                  Apresentação do projeto na integração dos alunos ingressantes e aos demais alunos e professores. \\ \hline 
 Reunião do P.I&                            03/08/2026&                03/08/2026&                 Tiago Rodrigues; Davi Almeida; Leandro Augusto;&                                 Reunião realizada para organizar as ideias do projeto, definir novas tarefas para o semestre e separar funções entre os integrantes.\\ \hline
                  Definição de metas para o projeto e cronograma.&                            03/08/2026&                10/08/2026&                 Tiago Rodrigues; Davi Almeida; Leandro Augusto;&                                  Definição de metas para o projeto e cronograma com auxilio do professor Yuri.\\ \hline
                 
\end{tabular}


\end{table}

 \begin{table}[]
\centering


\begin{tabular}{|>{\raggedright\arraybackslash}p{0.2\linewidth}|l|l|>{\raggedright\arraybackslash}p{0.2\linewidth}|>{\raggedright\arraybackslash}p{0.2\linewidth}|}

\hline
 Reformulação da paleta de cor&                            17/08/2026&                18/08/2026&                 Tiago Rodrigues; Leandro Augusto;&                                 Reformulação da paleta de cor do projeto para ficar alinhado com o atual design da aplicação.\\ \hline
                  Alterações no artigo&                            18/08/2026&               29/09/2026&                 Tiago Rodrigues.&                                  Alterações no artigo cientifico para corrigir pontos importantes e organizar ideias.\\ \hline
 
                  Modificação no arduíno&                            20/08/2026&                15/09/2026&                  Leandro Augusto;&                                 Alterações no código do arduíno e no modelo 3D do projeto para melhorar a performance e juntar com o alimentador.\\ \hline
                  Modificação no arduíno&                            29/08/2025&                29/10/2025&                  Leandro Augusto;&                                 Alterações no código do arduíno e no modelo 3D do projeto para melhorar a performance e ser de fácil uso.\\ \hline
 
\end{tabular}


\end{table}

 \begin{table}[]
\centering


\begin{tabular}{|>{\raggedright\arraybackslash}p{0.2\linewidth}|l|l|>{\raggedright\arraybackslash}p{0.2\linewidth}|>{\raggedright\arraybackslash}p{0.2\linewidth}|}

\hline

Visita na UNESP Registro para pesquisa de campo& 15/09/2025& 15/09/2025&  Tiago Rodrigues; Victor Roder; Leandro Augusto. &Realizada a primeira pesquisa de campo na UNESP Registro. Foi conduzida pela professora Giovana Bertini e pela aluna Estephany Konesuk\\\hline
 Reunião do P.I& 17/09/2025& 17/09/2025& João Kusaka; Matheus Abrahão; Tiago Rodrigues; Victor Roder; Leandro Augusto;& Realizada uma reunião na FATEC para adição de uma nova funcionalidade no sistema, para organizar pontos importantes e sobre o andamento do projeto.\\\hline
 Inclusão da funcionalidade de tutorial& 18/09/2025& 20/09/2025& João Kusaka.&Realizada a inclusão da funcionalidade de tutorial para facilitar a utilização para o usuário. A funcionalidade foi decidida na aula de Experiência de Usuário.\\ \hline 
 Coleta de camarões& 22/09/2025& 22/09/2025& Leandro Muniz. &Feita a coleta dos camarões que foram doados para o projeto pela professora Giovana Bertini da UNESP Registro\\\hline
                  Refazimento dos diagramas&                            29/09/2025&                29/10/2025&                 Tiago Rodrigues.&                                  Realizada as alterações nos diagramas do projeto com base em todas as mudanças realizadas no 4° semestre\\ \hline 

\end{tabular}


\end{table}

 \begin{table}[]
\centering


\begin{tabular}{|>{\raggedright\arraybackslash}p{0.2\linewidth}|l|l|>{\raggedright\arraybackslash}p{0.2\linewidth}|>{\raggedright\arraybackslash}p{0.2\linewidth}|}

\hline

 Workshop sobre o artigo& 10/10/2025& 10/10/2025& Tiago Rodrigues.&Workshop realizado pelo professor Frederico para auxiliar nas melhoras dos artigos e dicas.\\\hline
                  Modificações no Banner&                            13/10/2025&                15/10/2025&                 Tiago Rodrigues.&                                  Foram realizadas mudanças no banner para estarem com as informações atualizadas acerca do projeto e dos novos protótipos de tela. Banner atualizado pelo Figma.\\ \hline
                  Modificações no Pitch&                            16/10/2025&                29/10/2025&                 João Kusaka; Leandro Augusto; Matheus Ferrari; Tiago Rodrigues&                                  Foram realizadas alterações no pitch com base nos requisitos do MindMap\\ \hline                
\end{tabular}


\end{table}

\end{document}